\section{Introduction}
\IEEEPARstart{E}{ncrypted} communication hides payload content but leaves observable packet sizes, directions, and timing. Website fingerprinting demonstrates that these observations can support inference despite encryption~\cite{df,automated}. Flow-correlation research further shows that relationships between observations can be learned across a network path~\cite{deepcorr}. These findings motivate a different question for proxy traffic: can knowledge learned at one side of a proxy be reused when only the other side is available at inference time?

This question is not answered by high classification accuracy on either side alone. A model may distinguish deployments without identifying business activity. A feature may be accurately recoverable across sides while adding little to a particular classifier. Conversely, a mapping with imperfect reconstruction may preserve the relative class scores that matter to that classifier. Treating these properties as interchangeable risks attributing prediction improvements to a general invariant representation that has not been demonstrated.

Two challenges make the distinction operationally important. First, paired traffic shares workload, content, and visit-specific variation. Exact correspondence may expose any combination of these components, whereas a downstream label may distinguish only domain--activity classes. A correspondence objective can therefore favor faithfully preserving variation that is not essential to that classification problem. Second, the proxy entrance is not an environment-independent reference. Network feedback, implementations, and collection conditions can affect entrance traffic as well as exit traffic. Even perfect reconstruction of a target-deployment entrance representation does not guarantee that a source-deployment classifier will classify it correctly.

We address these challenges through a measurement-driven study with a fixed-capacity adaptation mechanism. We first characterize paired statistical changes and contrast feature recoverability with classification utility. We then retain a linear source classifier and compare direct application, coordinate-wise statistical calibration, and a calibration objective that also preserves centered source logits. Exact visit pairing, within-content cyclic mismatch, and content--deployment center targets isolate different uses of correspondence. Finally, we refit the complete pipeline on one deployment and evaluate held-out contents in the other deployment. This last experiment separates fit-time isolation from the historical fact that both deployments have already informed the research process.

The results show why a single ``pairing helps'' statement is insufficient. Under the same content partition, direct cross-side application reduces macro-F1 from \FOneA{} to \FOneB{}. Statistical calibration restores \FOneC{}, and decision-aware calibration reaches \FOneF{} with the same 298 mapping parameters. However, the decision-aware mismatched and center controls achieve \FOneG{} and \FOneH{}, respectively; the exact-target variant is not uniformly superior. Cross-deployment results are directional: calibration improves classification in the VLESS-to-SS direction, whereas the reverse direction chiefly improves probability loss. The target-entrance diagnostic also degrades in both directions.

Our contributions are threefold.
\begin{itemize}
\item We provide an auditable paired-traffic characterization that distinguishes repeatable statistical correspondence from finite-learner task utility, without using endpoint identity fields as predictor inputs.
\item We construct a fixed-source, fixed-capacity calibration study that separates recovery objective from correspondence granularity, retaining negative controls and deployment-specific exceptions.
\item We identify distinct cross-side and cross-deployment limits of model reuse through fit-time isolated evaluation and explicit observation permissions.
\end{itemize}
The contribution is the controlled empirical relationship among these properties. We do not claim a new information-theoretic identity, state-of-the-art classification, or an end-to-end operational attack.
